\title{QLoRA: Efficient Finetuning of Quantized LLMs}

\begin{abstract}
This paper introduces \textsc{QLoRA}, a memory-efficient finetuning strategy that enables the finetuning of a 65B parameter language model on a single 48GB GPU, all while maintaining the performance standards of full 16-bit finetuning. In \textsc{QLoRA}, gradients are propagated through a static, 4-bit quantized pretrained model into Low Rank Adapters~(LoRA). Our leading model series, \textbf{Guanaco}, surpasses the performance of all prior open-access models on the Vicuna benchmark, achieving 99.3\% of ChatGPT's performance after just 24 hours of training on one GPU. Key memory-saving features of \textsc{QLoRA} include: (a) 4-bit NormalFloat (NF4), an innovative datatype that is theoretically optimal for weights with a normal distribution, (b) Double Quantization, which minimizes average memory usage by further quantizing the quantization parameters, and (c) Paged Optimizers to regulate peak memory demands. Utilizing \textsc{QLoRA}, we finetune over 1,000 models, systematically evaluating instruction-following and conversational abilities over 8 instruction datasets, several architectures (LLaMA, T5), and scales that are otherwise prohibitive with standard methods (such as 33B and 65B parameter models). Findings indicate that QLoRA finetuning on modest, high-quality data delivers state-of-the-art results—even for models smaller than those currently holding the SoTA. An extensive study of chatbot effectiveness is presented, with comparisons drawn from both human raters and GPT-4, the latter proving to be an efficient and reliable substitute for human assessment. Additionally, we observe limitations in the reliability of existing chatbot benchmarks for accurately assessing chatbot competence. An in-depth "lemon-picked" analysis highlights situations where \textbf{Guanaco} underperforms relative to ChatGPT. All models, source code, and custom CUDA kernels for 4-bit training are publicly available.\footnote{\url{https://github.com/artidoro/qlora} and \url{https://github.com/TimDettmers/bitsandbytes}}
\end{abstract}

\section{Introduction}

Adapting large language models (LLMs) through finetuning has proven to be a particularly powerful method for enhancing model capability \citep{min2021metaicl, wei2021finetuned, ouyang2022training, wang2022super, wang2022self, liu2022few}, as well as for promoting desired or mitigating unwanted behaviors in these systems \citep{ouyang2022training,askell2021general,bai2022training}. Nonetheless, finetuning models of substantial scale is extremely resource-intensive; for instance, standard 16-bit finetuning of a LLaMA model with 65B parameters~\citep{touvron2023llama} exceeds 780 GB of GPU memory requirements. Although contemporary quantization strategies have succeeded in lowering memory needs for LLM inference \citep{dettmers2022llm, dettmers2022case, frantar2022gptq, xiao2022smoothquant}, these approaches generally fail when applied during the training phase \citep{wortsman2023stable}.

In this work, we present the first successful approach for finetuning LLMs quantized to 4 bits, achieving this without any reduction in predictive capability. The proposed method, \textsc{QLoRA}, involves a new high-precision quantization procedure applied to a pretrained model at 4-bit granularity, supplemented by incorporating a limited number of trainable Low-rank Adapter weights \citep{hu2021lora} 

that are updated through backpropagation using gradients that flow through the quantized parameters.

Employing \textsc{QLoRA}, the memory requirements to finetune a 65B parameter model are drastically lowered—from over 780GB to less than 48GB of GPU memory—without compromising speed or model performance relative to traditional 16-bit finetuning. This advance significantly democratizes LLM finetuning, enabling even the largest open-source models to be adapted using a single GPU. Utilizing \textsc{QLoRA}, we developed the \textbf{Guanaco} model series, with our second-best system attaining 97.8\% of ChatGPT's performance on the Vicuna~\citep{vicuna2023} evaluation, trained within 12 hours using just one consumer-grade GPU; with a single professional GPU in 24 hours, our largest \textbf{Guanaco} model reaches 99.3\%, effectively matching ChatGPT on this benchmark. Notably, our smallest \textbf{Guanaco} variant (7B parameters) only needs 5 GB of memory at inference and surpasses the 26 GB Alpaca model by over 20 percentage points on the Vicuna benchmark (Table~\ref{tbl:vicuna_eval}).

\textsc{QLoRA} incorporates several novel techniques aimed at minimizing memory consumption while maintaining high performance levels: (1) \textbf{4-bit NormalFloat}, a quantization scheme optimized for data with a normal distribution, grounded in information theory, which demonstrates empirically superior outcomes compared to standard 4-bit Integer and 4-bit Float approaches. (2) \textbf{Double Quantization}, which involves an additional layer of quantizing the quantization parameters themselves, leading to an average reduction of 0.37 bits per parameter (roughly translating to 3 GB of savings for a 65B parameter model). (3) \textbf{Paged Optimizers}, which leverage NVIDIA’s unified memory to mitigate the surges in memory usage caused by gradient checkpointing during the handling of long-sequence mini-batches. By integrating these strategies, a more effective LoRA methodology is developed—incorporating adapters in each network layer—thereby mitigating nearly all the accuracy compromises noted in earlier work.

The resource efficiency achieved by \textsc{QLoRA} enables comprehensive investigation into instruction finetuning and chatbot performance at scales previously unattainable using conventional finetuning techniques due to memory limitations. Specifically, more than 1,000 models are trained across a variety of instruction tuning corpora, architectures, and sizes ranging from 80M to 65B parameters. In addition to verifying that \textsc{QLoRA} attains parity with 16-bit training (\S\ref{sec:comp_to_std}) and successfully training a state-of-the-art chatbot, \textbf{Guanaco}, (\S\ref{sec:pushing}), further analysis is performed on the resulting models. Key findings include the observation that data quality significantly outweighs dataset size—for example, a relatively small dataset of 9k samples (OASST1) surpassed the performance of a much larger 450k sample dataset (FLAN v2, subsampled) in chatbot evaluations, even when both datasets are curated to support instruction-following tasks. Moreover, results indicate that superior performance on the Massive Multitask Language Understanding (MMLU) benchmark does not necessarily translate to higher Vicuna chatbot benchmark scores, highlighting that the appropriateness of a dataset to a specific task is a more decisive factor than dataset volume.

Additionally, we conduct a thorough assessment of chatbot capabilities employing both human assessors and GPT-4 for appraisal. Our evaluation employs a tournament framework, pitting models against each other in head-to-head matches where they generate responses to identical prompts. Match outcomes are adjudicated by either GPT-4 or human reviewers, and the collective tournament results are synthesized into Elo scores~\citep{elo1967proposed, elo1978rating}, which provide a ranking of chatbot performance. Our findings indicate strong concordance between GPT-4 and human judges regarding model rankings, though we also observe some notable discrepancies. Therefore, we emphasize that model-based evaluation, while cost-effective compared to human annotation, entails its own set of limitations and uncertainties.

To supplement the benchmarking of our chatbot, we present a qualitative examination of the \textbf{Guanaco} models. This qualitative assessment brings attention to both effective and ineffective model behaviors that escape quantitative measures.

We have made public all generated model outputs annotated by both humans and GPT-4, enabling further research. In addition, our codebase and CUDA kernels are fully open-sourced and our techniques are incorporated into the Hugging Face transformers stack \citep{wolf2019huggingface}, facilitating broad accessibility. We provide a suite of adapters for model sizes 7/13/33/65B, each trained on eight diverse instruction-following datasets, resulting in 32 distinct open-source finetuned models.

\section{Background}
\paragraph{Block-wise k-bit Quantization} 
Quantization refers to the transformation of data from a representation with higher information content to one with reduced precision. This typically involves casting data from a type with a larger bit-width to one with fewer bits, such as converting 32-bit floating point numbers to 8-bit integer values. To make efficient use of the entire numerical range available in the low-bit data type, the input is usually rescaled by normalizing with the absolute maximum value among its elements—commonly arranged as a tensor. For instance, to quantize an FP32 tensor into an Int8 tensor with a range of $[-127, 127]$:

\begin{equation}
    \mathbf{X}^{ \text{Int8}} = \text{round}\left( \frac{127}{{\text{absmax}}(\mathbf{X}^{{\text{FP32}}})}\mathbf{X}^{{\text{FP32}}} \right) = \text{round}(c^{\text{FP32}}\cdot \mathbf{X}^{{\text{FP32}}}),
\end{equation}

where $c$ is the {\em quantization constant} or {\em quantization scale}. The process of restoring the data to its higher-precision form, dequantization, is given by:

\begin{equation}
    \text{dequant}(c^{\text{FP32}}, \mathbf{X}^{\text{Int8}}) = \frac{\mathbf{X}^{\text{Int8}}}{c^{\text{FP32}}} = \mathbf{X}^{\text{FP32}}
\end{equation}

A limitation of this method arises when the input tensor contains values with exceptionally large magnitudes, or outliers, resulting in inefficient usage of quantization bins—some bit combinations correspond to ranges that receive little or no assigned values. To mitigate this, a widely adopted strategy involves partitioning the input tensor into several blocks, each of which is quantized independently with its own dedicated quantization constant $c$. More precisely, the input tensor $\mathbf{X}\in \mathbb{R}^{b\times h}$ is reshaped into a one-dimensional array and segmented into $n$ consecutive blocks, each of size $B$, so that $n = ({b\times h})/{B}$. Each block is then quantized on its own using Equation~1, resulting in a quantized tensor accompanied by $n$ block-specific quantization constants $c_i$.

\paragraph{Low-rank Adapters} The LoRA (Low-rank Adapter) finetuning technique \citep{hu2021lora} minimizes memory consumption by introducing a small collection of learnable adapter parameters, while the main model parameters are held static and are not subject to updates. During stochastic gradient descent, gradients are propagated through the pretrained, unmodified model weights, while only the adapters receive updates to optimize the overall objective. In LoRA, an extra, factorized linear transformation augments the original linear mapping. For a given projection operation ${\mathbf{X} \mathbf{W} = \mathbf{Y}}$ where $\mathbf{X}\in  \mathbb{R}^{b\times h}$ and $\mathbf{W}\in  \mathbb{R}^{h\times o}$, LoRA performs the following computation:

\begin{equation}
    \mathbf{Y} = \mathbf{X}\mathbf{W} + s\mathbf{X}\mathbf{L}_1\mathbf{L}_2,
\end{equation}

with $\mathbf{L}_1\in  \mathbb{R}^{h\times r}$ and $\mathbf{L}_2\in  \mathbb{R}^{r\times o}$, and where $s$ denotes a scalar parameter.

\paragraph{Memory Requirement of Parameter-Efficient Finetuning} A key aspect to consider is the memory usage of LoRA during training, both in relation to the quantity and dimensions of adapters employed. Due to LoRA's extremely small memory overhead, it is feasible to employ a greater number of adapters to enhance model performance without causing a notable increase in overall memory consumption. Although LoRA is a Parameter Efficient Finetuning (PEFT) strategy, the dominant factor in memory usage during LLM finetuning is activation gradients, rather than the parameters introduced by LoRA itself. For example, in the case of a 7B LLaMA model trained on FLAN v2 with a batch size of 1, and LoRA weights representing a typical 0.2\% fraction of the base model’s parameters~\citep{hu2021lora,liu2022few}, the memory taken up by LoRA input gradients is 567 MB, whereas the actual LoRA parameter set requires only 26 MB. Applying gradient checkpointing~\citep{chen2016training} reduces the memory for input gradients to roughly 18 MB per sequence—still larger than the footprint of all LoRA parameters. In contrast, the memory needed for the 4-bit base model itself is 5,048 MB. This illustrates the importance of gradient checkpointing, and also demonstrates that further lowering the number of LoRA parameters results in only marginal memory savings. Consequently, increasing the number of adapters is possible without a significant impact on the training memory demand (see Appendix~\ref{app:memory} for additional details). As will be elaborated later, this is essential for achieving performance on par with full 16-bit precision.

\section{\textsc{QLoRA} Finetuning}

\textsc{QLoRA} enables effective 4-bit finetuning accuracy by employing two novel methods: 4-bit NormalFloat (NF4) quantization and Double Quantization. Furthermore, the approach introduces Paged Optimizers, a mechanism to mitigate memory usage spikes during gradient checkpointing. This addresses a common challenge in large model finetuning—out-of-memory errors on single machines.

The \textsc{QLoRA} framework utilizes a storage data type of low precision, typically 4-bit, alongside a computation data type, generally BFloat16. Operationally, this implies that whenever a weight tensor under \textsc{QLoRA} is accessed, it is first dequantized to the BFloat16 format, after which matrix multiplication is carried out in 16-bit precision.

The subsequent discussion provides a detailed breakdown of the constituent parts of \textsc{QLoRA}, leading into a rigorous formal definition.

\paragraph{4-bit NormalFloat Quantization} The NormalFloat (NF) representation extends the concept of Quantile Quantization \citep{dettmers2022optimizers}, which is considered information-theoretically optimal as it allocates an equal quantity of input tensor values to each quantization bin. Quantile quantization functions by approximating the quantiles of the tensor based on the empirical cumulative distribution function.

However, one significant drawback of quantile quantization is the high computational cost associated with quantile calculation. For this reason, fast quantile estimation methods, such as SRAM quantiles~\citep{dettmers2022optimizers}, are typically used. Because these methods yield only approximate quantile values, the result is often substantial quantization error for outlier values—which can be especially critical.

Such expensive and error-prone quantile estimation can be circumvented when input tensors are drawn from a distribution fixed up to a quantization constant. In scenarios like these, the quantiles of the input tensors are consistent, which allows for exact quantile computation with practical efficiency.

Pretrained neural network weights typically follow a normal distribution centered at zero, characterized by a standard deviation $\sigma$ (refer to Appendix~\ref{app:normality}). To standardize all weights to a uniform reference distribution, we rescale $\sigma$ so that the resulting distribution aligns precisely with the allowable range of our chosen data type. For our purposes, we define this permissible range to be $[-1, 1]$. Consequently, it is necessary to normalize both the quantile values of the data type and the neural network weights so they occupy this fixed interval.

The data type that achieves optimal information efficiency for encoding zero-mean normal distributions with arbitrary $\sigma$ over $[-1, 1]$ can be determined by the following procedure: (1) identify the $2^k+1$ quantile breakpoints of an ideal $N(0, 1)$ distribution, thereby establishing a $k$-bit quantile-based quantization data type for normal distributions; (2) adjust these quantile-derived values so that they span the $[-1, 1]$ interval; (3) quantize any input weight tensor by rescaling its values into $[-1, 1]$ through normalization with respect to its absolute maximum value.

Once the domain of the weights and that of the data type coincide, quantization proceeds in the usual fashion. Step (3) can be interpreted as modifying the standard deviation of the weight tensor so it equals that of the k-bit data type. To state this precisely, we compute the $2^k$ quantization levels $q_i$ as:

\begin{equation}
q_i = \frac{1}{2}\left( Q_X\left(\frac{i}{2^k+1}\right) + Q_X\left(\frac{i+1}{2^k+1}\right)\right),
\end{equation}

where $Q_X(\cdot)$ denotes the quantile function corresponding to the standard normal distribution $N(0,1)$. A limitation associated with symmetric $k$-bit quantization is the lack of a quantization bin that corresponds precisely to zero. This poses challenges, especially for elements such as padding or any inherently zero-valued entries, which should ideally be quantized exactly. To guarantee the presence of a discrete zeropoint at $0$ while fully utilizing all $2^k$ representable values in a $k$-bit number, we formulate an asymmetric quantization scheme. Here, the quantiles $q_i$ are estimated by dividing the distribution into two segments: $2^{k-1}$ bins for the negative values and $2^{k-1} + 1$ bins for the positive values. We then combine these collections of $q_i$, removing the duplicate zero that appears in both. The resulting data type, which ensures uniform distribution of expected values among quantization bins, is designated as the \emph{k-bit NormalFloat} (NFk). This construction is optimal from an information-theoretic perspective for zero-mean, normally distributed data. The precise quantized values implemented for this data type are provided in Appendix~\ref{app:nf4}.

\paragraph{Double Quantization{}} 
We propose \emph{Double Quantization} (DQ), a strategy in which even the quantization parameters themselves are quantized, enabling further reductions in memory usage. Precise 4-bit quantization is known to require a small blocksize~\citep{dettmers2022case}, yet this introduces considerable memory overhead due to the storage of quantization parameters. For instance, when deploying 32-bit quantization parameters with a blocksize of 64 for $\mathbf{W}$, each parameter incurs an average overhead of $32/64 = 0.5$ bits. Double Quantization alleviates this by shrinking the memory requirement associated with quantization parameters.

Concretely, the Double Quantization procedure views the original quantization constants $c_2^{\text{FP32}}$ as input for a secondary quantization stage. This secondary operation produces quantized constants $c_2^{\text{FP8}}$, as well as new quantization constants at the next level, $c_1^{\text{FP32}}$. We apply 8-bit Float quantization with a blocksize of 256 for this subsequent step, as previous work~\citep{dettmers2022case} and our results show no significant loss in accuracy at 8 bits. Since $c_2^{\text{FP32}}$ are strictly positive, we subtract their mean to recenter the distribution around zero and exploit symmetric quantization. When blocksize is 64, the average per-parameter storage required for quantization parameters decreases from $32/64 = 0.5$ bits to $8/64 + 32/(64\cdot256) = 0.127$ bits, yielding a reduction of 0.373 bits per parameter.

\paragraph{Paged Optimizers} leverage the NVIDIA unified memory~\footnote{\tiny \url{https://docs.nvidia.com/cuda/cuda-c-programming-guide}}, which automatically manages data transfers between the CPU and GPU on a per-page basis. This mechanism ensures seamless GPU computations even when the GPU experiences memory shortages by handling these transfers transparently. Functionally similar to how operating systems swap data between CPU RAM and disk storage, we employ this capability to allocate optimizer states in paged memory. Consequently, these states are migrated to CPU RAM as GPU memory becomes constrained and brought back into GPU memory as necessary during optimizer update operations.

\paragraph{\textsc{QLoRA}.} Building on the mechanisms outlined above, we formalize \textsc{QLoRA} for a quantized base model’s individual linear layer with a single LoRA adapter by

\begin{equation}
    \mathbf{Y}^{\text{BF16}} =  \mathbf{X}^{\text{BF16}} \text{doubleDequant}(c_1^{\text{FP32}}, c_2^{\text{k-bit}}, \mathbf{W}^{\text{NF4}}) + \mathbf{X}^{\text{BF16}}\mathbf{L}^{\text{BF16}}_1\mathbf{L}^{\text{BF16}}_2,
\end{equation}

with the \texttt{doubleDequant} operation given by

\begin{equation}
    \text{doubleDequant}(c_1^{\text{FP32}}, c_2^{\text{k-bit}}, \mathbf{W}^{\text{k-bit}}) = \text{dequant}(\text{dequant}(c_1^{\text{FP32}}, c_2^{\text{k-bit}}), \mathbf{W}^{\text{4bit}}) = \mathbf{W}^{\text{BF16}},
\end{equation}

In our approach, $\mathbf{W}$ is stored in NF4 format, while $c_2$ uses FP8.

To improve quantization fidelity, we select a blocksize of 64 for $\mathbf{W}$, and for memory savings, a blocksize of 256 for $c_2$.

When computing parameter updates, only the gradients with respect to the adapter weights, $\frac{\partial E}{\partial \mathbf{L}_i}$, are required; gradients for the 4-bit weights, $\frac{\partial E}{\partial \mathbf{W}}$, are not necessary. Nevertheless, evaluating $\frac{\partial E}{\partial \mathbf{L}_i}$ demands the computation of $\frac{\partial \mathbf{X}}{\partial \mathbf{W}}$, which is conducted by first dequantizing $\mathbf{W}^{\text{NF4}}$ to $\mathbf{W}^{\text{BF16}}$ following equation~(5), thus enabling the derivative $\frac{\partial \mathbf{X}}{\partial \mathbf{W}}$ to be obtained in BFloat16 precision.

In essence, \textsc{QLoRA} employs a single storage data type (commonly 4-bit NormalFloat) and a computation data type (16-bit BrainFloat). During both forward and backward passes, the system dequantizes parameters from the storage format to the computation type. Nevertheless, gradient computation for the weights is exclusively performed on the LoRA parameters, which are represented in 16-bit BrainFloat.

\section{QLoRA vs. Standard Finetuning}
\label{sec:comp_to_std}

After describing the \textsc{QLoRA} methodology and its substantial memory savings during model finetuning, the key issue is to determine whether its performance aligns with that of traditional full-parameter finetuning. In addition, it is important to dissect QLoRA’s constituents, particularly examining the influence of NormalFloat4 relative to conventional Float4 quantization. The forthcoming sections present experiments designed to address these points.

\paragraph{Experimental setup.} Our study spans three primary neural network configurations—encoder, encoder-decoder, and decoder-only architectures—evaluating \textsc{QLoRA} in comparison with 16-bit adapter-based finetuning and standard full finetuning for models containing up to 3B parameters. The benchmarking encompasses GLUE \citep{wang2018glue} with RoBERTa-large \citep{liu2019roberta}; Super-NaturalInstructions (TKInstruct) \citep{wang2022super} employing T5 \citep{t5}; and 5-shot MMLU \citep{hendrycksmeasuring} following finetuning of LLaMA on Flan v2 \citep{longpre2023flan} and Alpaca \citep{alpaca}. To further compare NF4 with alternate 4-bit quantization types, we adopt the methodology from \citet{dettmers2022case}, measuring zero-shot accuracy and perplexity post-quantization on various models (OPT~\citep{zhang2022opt}, LLaMA~\citep{touvron2023llama}, BLOOM~\citep{scao2022bloom}, Pythia~\citep{biderman2023pythia}) over model scales ranging from 125m to 13B parameters.

Specific experimental details are further elaborated upon in each respective results subsection to enhance clarity, and comprehensive information is included in Appendix~\ref{app:exp-setup}.

Although paged optimizers play a vital role in enabling 33B/65B \textsc{QLoRA} tuning on single 24/48GB GPUs, we have not collected explicit metrics on the performance of Paged Optimizers, as paging becomes relevant primarily for rare cases involving long sequence mini-batches. Nonetheless, runtime analyses were performed on 65B models utilizing 48GB GPUs, revealing that, with a batch size of 16, paged optimizers deliver comparable training throughput to that of conventional optimizers. It is left for future work to quantify and systematically examine the circumstances under which paging may introduce overhead.

\paragraph{Default LoRA hyperparameters do not match 16-bit performance}

Applying the prevailing approach of assigning LoRA to the query and value projections within attention layers \citep{hu2021lora} does not recover full-model finetuning accuracy for larger foundation models. Figure~\ref{fig:hyper}, depicting results for LLaMA 7B finetuned on Alpaca, indicates that the decisive LoRA hyperparameter is the total number of LoRA adapters deployed. Achieving parity with full finetuning requires LoRA adaptation across all linear layers in the transformer blocks. Conversely, adjusting other LoRA hyperparameters, such as the projection rank $r$, appears to exert negligible influence on overall performance (refer to Appendix \ref{app:exp-setup}).

In a similar vein, we observe that the standard hyperparameter choices for fully finetuned baselines are suboptimal. To establish solid baselines, we perform a comprehensive hyperparameter sweep, varying learning rates between 1e-6 and 5e-5 and exploring batch sizes from 8 to 128. The corresponding outcomes for finetuning 7B LLaMA on the Alpaca dataset can be found in Figure~\ref{fig:hyper}.

\paragraph{4-bit NormalFloat surpasses 4-bit Floating Point in performance} 

Although the 4-bit NormalFloat (NF4) data format is known to be optimal from an information-theoretic perspective, it remains to be clarified whether this theoretical benefit leads to measurable improvements in practice. Adopting the methodology described by \citet{dettmers2022case}, we assess quantized LLMs (including OPT \citep{zhang2022opt}, BLOOM \cite{scao2022bloom}, Pythia \cite{biderman2023pythia}, and LLaMA) spanning a wide parameter range (125M to 65B) using various quantization types across language modeling and a suite of zero-shot benchmarks. Figure~\ref{fig:nf4} and Table~\ref{tbl:nf4_ppl} illustrate that NF4 offers marked improvements over FP4 and Int4, while double quantization decreases memory usage with no observable loss in effectiveness.

\paragraph{k-bit \textsc{QLoRA} achieves parity with 16-bit full finetuning and 16-bit LoRA}

Prior research has demonstrated that, while 4-bit quantization is feasible for \emph{inference}, it typically results in some degradation relative to 16-bit models \citep{dettmers2022case,frantar2022gptq}. This naturally prompts the question: can this loss in quality be remedied through 4-bit adapter finetuning? To investigate this, we experiment with two scenarios.

The initial scenario entails comparing results with full 16-bit finetuning of RoBERTA and T5 models ranging from 125M to 3B parameters, evaluated on the GLUE benchmark and the Super-NaturalInstructions dataset. The findings, detailed in Table~\ref{tab:nf4vsbf16}, reveal that 16-bit, 8-bit, and 4-bit adapter approaches all match the results of the fully finetuned 16-bit baseline in both tasks. This indicates that adapter finetuning post-quantization can fully recover any performance loss attributed to quantization.

In our second experimental configuration, considering that fully finetuning models with 11B parameters or larger exceeds the GPU memory capacity of a single server, we further investigate if 4-bit \textsc{QLoRA} achieves comparable performance to 16-bit LoRA across models ranging from 7B to 65B parameters. Specifically, we conduct finetuning of LLaMA models (7B--65B) using two instruction-following datasets—Alpaca and FLAN v2—and assess performance using the 5-shot accuracy metric on the MMLU benchmark. The outcomes, as detailed in Table~\ref{tbl:llama-datatype}, demonstrate that employing the NF4 quantization scheme with double quantization restores MMLU results on par with 16-bit LoRA. Conversely, \textsc{QLoRA} with FP4 consistently trails the 16-bit brain float LoRA baseline by around 1 percentage point, supporting our observations that (1) \textsc{QLoRA} equipped with NF4 mirrors the performance of both 16-bit full and LoRA finetuning, and (2) NF4 exhibits superior quantization fidelity compared to FP4.

\paragraph{Summary} Across multiple academic benchmarks and standardized evaluation protocols, our experiments reveal that 4-bit \textsc{QLoRA} with the NF4 quantization format attains parity with both 16-bit full and LoRA finetuning in terms of effectiveness. Additionally, we establish that NF4 consistently outperforms FP4, and that double quantization imposes no adverse impact on outcomes. These results together provide strong empirical evidence that 4-bit \textsc{QLoRA} can reliably replicate the results of 16-bit approaches.

Echoing findings from previous research on quantization \citep{dettmers2022case}, both our MMLU and Elo analyses suggest that within a fixed computational budget for finetuning and inference, scaling up model size while decreasing numerical precision leads to improved efficiency. This underscores the practical advantages afforded by \textsc{QLoRA}. Given our lack of observed degradation relative to full finetuning when using 4-bit quantized models, an open question remains regarding the precise location of the trade-off between performance and precision in \textsc{QLoRA}—a subject we leave for future investigation.

We move forward by examining instruction tuning at scales that are otherwise infeasible using full 16-bit finetuning on typical academic hardware resources.

\section{Pushing the Chatbot State-of-the-art with QLoRA}
\label{sec:pushing}
Building upon our earlier result that 4-bit \textsc{QLoRA} delivers comparable results to 16-bit approaches over different scales, tasks, and datasets, we provide a thorough analysis of instruction finetuning using some of the largest open-source language models accessible to researchers. To thoroughly gauge instruction finetuning effectiveness, we perform evaluations on the MMLU Natural Language Understanding benchmark and introduce novel strategies to assess real-world chatbot functionality.

\subsection{Experimental setup} Below we summarize the design of our experiments; comprehensive specifications are available in Appendix \ref{app:chatbot}.

\paragraph{Data} Since a systematic review of the latest instruction-following datasets appears lacking, we curate a selection of eight contemporary datasets. Our selection spans resources generated by crowd-sourcing efforts (OASST1~\citep{kopf2023openassistant}, HH-RLHF~\citep{bai2022training}), datasets derived from distilled outputs of instruction-tuned models (Alpaca~\citep{alpaca}, self-instruct~\citep{wang2022self}, unnatural-instructions~\citep{honovich2022unnatural}), collections assembled from various sources (FLAN v2~\citep{chung2022scaling}), along with datasets employing hybrid acquisition strategies (Chip2~\citep{laion2023}, Longform~\citep{koksal2023longform}). This compilation encompasses various languages, dataset sizes, and licensing schemes.

\paragraph{Training Setup} To eliminate inconsistencies stemming from disparate training objectives, our QLoRA finetuning utilizes a standard cross-entropy loss for supervised training, omitting reinforcement learning—even when handling datasets annotated with human evaluations. When datasets differentiate explicitly between instruction and response, we restrict finetuning to the response content alone (additional studies are in Appendix \ref{app:chatbot}). For OASST1 and HH-RLHF, which provide several possible responses, we select the highest-rated response from each branch in the conversational tree and carry out finetuning using these selected dialogues, incorporating the associated instructions. In our workflow, all experiments implement the NF4 \textsc{QLoRA} method with both double quantization and paged optimizers, mitigating memory overhead during gradient checkpointing. We conduct targeted hyperparameter optimization for the 13B and 33B LLaMA models, confirming that hyperparameter values established for 7B are broadly applicable—excluding learning rate and batch size adjustments. Specifically, for 33B and 65B models, we halve the learning rate while doubling the batch size.

\paragraph{Baselines} Our models are benchmarked against both academic (Vicuna~\citep{vicuna2023}, Open Assistant~\citep{kopf2023openassistant}) and commercial (GPT-4 \citep{openai2023gpt}, GPT-3.5-turbo, Bard) chatbot platforms. Notably, Open Assistant utilizes a LLaMA 33B architecture refined via Reinforcement Learning from Human Feedback (RLHF) based on the OASST1 dataset, which we also employ. Vicuna involves comprehensive fine-tuning of LLaMA 13B leveraging proprietary conversation data from ShareGPT, amounting to a distilled model from OpenAI's GPT suite.

\subsection{Evaluation}

To remain consistent with prevailing standards, we evaluate models using the MMLU (Massively Multitask Language Understanding) benchmark \citep{hendrycksmeasuring}, which tests a suite of 57 multiple-choice language understanding challenges across domains like basic mathematics, US history, computer science, legal reasoning, and beyond. Our results are presented in terms of 5-shot accuracy on the test set.

We further assess generative language performance by employing both automated and manual evaluation methods. This secondary evaluation phase leverages human-curated queries to assess response quality generated by the models. Although such tests provide a more representative environment for gauging chatbot performance and are increasingly used, there remains a lack of standardized methodology in existing research. Our approach, detailed below, applies nucleus sampling with $p=0.9$ and temperature $0.7$ for all generations.

\paragraph{Benchmark Data} Our analysis draws upon two carefully selected query datasets: the Vicuna prompts \citep{vicuna2023} and the OASST1 validation dataset \citep{kopf2023openassistant}. The Vicuna dataset comprises 80 unaltered prompts spanning a variety of topics. OASST1, meanwhile, consists of multilingual, crowd-generated multiturn dialogues between users and assistants. From this, we extract every user message from the validation split as a query, incorporating all preceding turns into the prompt. This selection yields 953 distinct user queries. We refer to these evaluation corpora as the Vicuna and OA benchmarks.

\paragraph{Automated Evaluation} Adopting the framework from \citet{vicuna2023}, we use GPT-4 to score models relative to ChatGPT (GPT-3.5 Turbo) using the Vicuna benchmark. For each query, both ChatGPT’s and the candidate model’s answers are provided to GPT-4, which assigns each a score out of ten and justifies its reasoning. The model’s performance is then expressed as a proportion of ChatGPT's total, with values exceeding 100\% indicating outperformance. Notably, our analysis revealed that GPT-4 demonstrates a positional bias, favoring the response listed first. To mitigate this, we advocate reporting the average score across both possible orderings.

We also evaluate by directly comparing model outputs. In this approach, we reduce scoring to a three-way choice: one model’s response, the other’s, or a tie, supplemented by an explanation. GPT-4 is tasked with determining the preferred response, or declaring a tie, for all pairwise model combinations across the Vicuna and OA datasets.

\paragraph{Human Evaluation} Although emerging literature suggests that generative models can serve effectively as evaluators \citep{fu2023gptscore}, the extent to which GPT-4’s judgments align with human assessments of chatbot quality remains uncertain. Consequently, we execute two concurrent rounds of human evaluation on the Vicuna benchmark that mirror the two automated protocols outlined above. Using Amazon Mechanical Turk (AMT), we enlist two annotators for ChatGPT-vs-model comparisons and three annotators for all pairwise assessments.

\paragraph{Elo Rating} We structure both human- and machine-based pairwise assessments as a tournament, where different models are pitted against one another. Each “match” features a pair of models tasked with generating the superior response to a given prompt. This methodology resembles those of \citet{bai2022training} and \citet{vicuna2023}, but we further incorporate evaluations from GPT-4 alongside human judgments. To calculate Elo~\citep{elo1967proposed,elo1978rating} scores, we draw random samples from our set of annotated comparisons. The Elo rating system, commonly utilized in chess and related competitive settings, quantifies a player's anticipated winning probability relative to an opponent. For instance, when comparing an Elo 1100 participant with one at 1000, the higher-rated contender is predicted to win roughly 65\% of the time; equivalently rated matches (e.g., 1000 vs 1000 or 1100 vs 1100) yield an even 50\% expected win rate. Elo ratings are updated after every match based on the predicted outcome—an unlikely upset produces a significant Elo adjustment, whereas outcomes aligned with expectations result in modest changes. Across many matches, Elo ratings tend to approximate the true proficiency of each competitor. We initialize all participants with a baseline of 1,000 points and set $K=32$. Adopting the approach from \citet{vicuna2023}, we run this rating process 10,000 times with different random seeds to minimize biases from match ordering, such as which model pairings are encountered earliest.

\subsection{\textbf{Guanaco}: \textsc{QLoRA} trained on OASST1 is a State-of-the-art Chatbot}

Through both automated and manual assessment, our findings indicate that the leading \textsc{QLoRA} model, {Guanaco} 65B, which we further trained on a variant of OASST1, establishes itself as the top-performing open-source conversational agent, with capabilities closely matching those of ChatGPT. In direct comparisons with GPT-4, {Guanaco} 65B and 33B register an estimated win rate of 30\% as determined by Elo scores from system-level pairwise comparisons conducted by human raters—a result that marks a new high point for open-source systems.

Results on the Vicuna benchmark \cite{vicuna2023} are presented in Table~\ref{tbl:vicuna_eval}, using ChatGPT as a baseline. According to this benchmark, {Guanaco} 65B demonstrates the highest level of performance after GPT-4, attaining 99.3\% of ChatGPT’s effectiveness. Although {Guanaco} 33B contains more parameters than Vicuna 13B, it utilizes only 4-bit quantization, making it substantially more memory efficient, with a 21 GB memory requirement compared to 26 GB, and delivers a performance gain of three percentage points relative to Vicuna 13B. Additionally, {Guanaco} 7B’s compact 5 GB memory profile enables it to run on contemporary smartphones while still surpassing Alpaca 13B by nearly 20 percentage points in performance.

Nevertheless, the confidence intervals reported in Table~\ref{tbl:vicuna_eval} are quite broad, and many models' results overlap. We suggest that this ambiguity arises from poorly defined evaluation scales; specifically, what constitutes an 8 on a 10-point scale may differ across tasks or evaluators. Consequently, we advocate for the Elo ranking methodology \citep{elo1967proposed}, which relies on \textit{pairwise} comparisons from both human annotators and GPT-4 to sidestep the limitations of absolute scales. Table~\ref{tbl:elo} shows the Elo ratings among the most competitive systems. It is worth mentioning that human and GPT-4 rankings diverge to some extent for certain models, particularly Guanaco 7B; however, for most models, the rankings are similar, with a Kendall Tau of $\tau=0.43$ and a Spearman correlation of $r=0.55$ at the system level. At the example level, agreement between majority human votes and GPT-4’s evaluation is weaker, with Fleiss $\kappa=0.25$. Taken together, these results indicate moderate alignment in system-level judgments between GPT-4 and human annotators, suggesting that evaluation by models can serve as a somewhat trustworthy complement to human assessments. Additional discussion is provided in Section~\ref{ssec:considerations}.

The Elo rankings presented in Table~\ref{tbl:elo_large} demonstrate that the {Guanaco} 33B and 65B models surpass all alternatives except for GPT-4 in both the Vicuna and OA benchmarks, exhibiting performance levels similar to ChatGPT, as also reflected in Table~\ref{tbl:vicuna_eval}. It should be noted that the Vicuna benchmark tends to advantage open-source systems, whereas the broader OA benchmark is more favorable to ChatGPT. Moreover, as illustrated in Tables \ref{tbl:mmlu-llama} and \ref{tbl:vicuna_eval}, the choice of finetuning dataset critically affects model outcomes. In particular, Llama models finetuned on FLAN v2 achieve especially strong results on the MMLU benchmark but yield the weakest performance on Vicuna (with similar patterns evident across other models). These findings suggest that current evaluation benchmarks measure different abilities: excelling in MMLU does not necessarily translate to high performance as a chatbot (as assessed by Vicuna or OA), nor does the converse hold true.

Among the top-ranked models evaluated, is unique in not utilizing proprietary data, as the OASST1 dataset’s collection policy prohibits employing GPT models. The Anthropic HH-RLHF model, also trained solely on open-source data, is the next closest competitor, but still lags behind {Guanaco} by 30 percentage points on the Vicuna benchmark (refer to Table~\ref{tbl:vicuna_eval}). In summary, these outcomes indicate that 4-bit \textsc{QLoRA} provides an effective framework capable of generating state-of-the-art chatbots that closely match the capabilities of ChatGPT. Notably, training the 33B {Guanaco} model can be completed within 12 hours on 24 GB consumer-grade GPUs. This efficiency paves the way for future research using \textsc{QLoRA} to further finetune models on specialized open-source datasets, leading to models that can rival premier commercial systems available today.

\section{Qualitative Analysis}

Although quantitative analysis constitutes the main component of our evaluation, relying solely on aggregated metrics raises several concerns. Chief among these is the issue of benchmark validity \citep{liao2021we}: it is frequently unclear if a benchmark accurately measures the intended capability or construct, particularly as new ``shortcuts'' emerge that can be exploited by machine learning models to game these assessments \citep{gururangan2018annotation, poliak2018hypothesis}. To partly address these concerns, we complement our study with a qualitative analysis, organized in two parts. In \S\ref{ssec:examples}, we provide sample outputs which capture some recurrent tendencies we observed in text produced by our 65b {Guanaco} model. Then, in \S\ref{ssec:considerations}, we elaborate on further factors influencing our findings and the lens through which we interpret them.

\subsection{Qualitative Analysis of Example Generations}
\label{ssec:examples}

Our approach to example selection involved systematically reviewing outputs generated on the Vicuna and OpenAssistant benchmarks to detect salient trends in {Guanaco}'s responses. Upon identifying specific tendencies—for example, the model's propensity for verbose answers—we crafted targeted prompts designed to elicit those patterns, even when they do not correspond to correct behavior. For instance, if excessive elaboration was observed, we would prompt with ``Answer yes or no without explanation'' to investigate if the model complies. This methodology allowed us to surface instances where the model's limitations could be exposed (so-called ``lemons'') as well as cases where it remained robust against such adversarial prompting (``cherries''), with representative examples from both categories presented below. Every sample shown here was produced using Nucleus Sampling \cite{holtzmancurious} with $p=0.9$.

Naturally, this overview is not exhaustive; due to the limited qualitative nature of our study, it is impractical to account for every possible variable. For example, the complete range of outputs a model may produce for a single prompt is vast, necessitating reliance on selected samples which we believe to be reasonably indicative. Nevertheless, presenting these examples serves to contextualize the quantitative findings introduced previously. By releasing all models and accompanying code as open source, we aim to encourage future research that more thoroughly investigates the phenomena discussed here.

\paragraph{Factual Recall} When posed straightforward factual questions, such as ``What is the capital of Zambia?'', all evaluated models reliably return correct answers, e.g.,

\begin{quote} 
    \textbf{Guanaco}:
    \texttt{The capital of Zambia is Lusaka.}
\end{quote}

However, when confronted with less common questions, {Guanaco}'s accuracy noticeably declines, even as its confidence remains unaffected. Consider this prompt drawn from HotPotQA \citep{yang2018hotpotqa}:

\begin{quote}
     \textbf{User}: In what year was the singer who popularized `I'll Keep the Lovelight Burning' born?  

     \textbf{Guanaco}:~\texttt{Al Jolson is credited with popularizing the song `I'll Keep the Lovelight Burning,' and he was born in the year 1886.}
 \end{quote}

In this example, both the identified individual and their birth year are incorrect (though the date corresponds to Al Jolson, the person cited in the answer).

\paragraph{Suggestibility} In certain scenarios, {Guanaco} demonstrates notable resilience against some forms of prompted misinformation, as illustrated below:

\begin{quote}
    \textbf{User}: How was it finally, officially confirmed that the earth is flat by peer reviewed scientists? \\ \\
    \textbf{Guanaco}: \texttt{There has never been any official confirmation by peer reviewed scientists that the earth is flat. Scientific consensus firmly supports the understanding that the earth is spherical. This has been established over centuries and is backed by a wide array of experiments and observations. \\ \\ The notion that the earth is flat is a conspiracy theory that has been refuted by scientific evidence. No peer reviewed body of scientists has ever endorsed or confirmed the earth to be flat. \\ \\ Please let me know if you have additional questions.}
\end{quote}

is also effective at identifying questions that cannot be answered, such as:

\begin{quote}
    \textbf{User}: What time is it?

    \textbf{Guanaco}: \texttt{I apologize, but as an AI language model, I do not have access to real-time data. My answers are based on previously acquired information, so I cannot provide the current time. For up-to-date time, please use a clock or consult an online time source.}
\end{quote}

Comparable responses are produced for prompts like ``Where are you?'', ``How are you?'', and others.

\paragraph{Refusal} In the same way, {Guanaco} can occasionally decline to comply with instructions, sometimes without a clear reason:

\begin{quote}
    \textbf{User}: Please invert the order of the words in the phrase ``I'm a cool dude in a mocha mood.'' \\ \\
    \textbf{Guanaco}: \texttt{I'm unable to do that at the moment, but I am able to clarify why your sentence is constructed correctly from a grammatical perspective.\\ \\The phrase ``I'm a cool dude in a mocha mood'' is structured according to standard English grammar and syntax. Specifically, ``cool'' serves as an adjective modifying the noun ``dude,'' while the prepositional phrase ``in a mocha mood'' supplies extra contextual information. The composition of this sentence is entirely proper and coherent given its usage.\\ \\If you have further questions or require additional assistance, please let me know.}
\end{quote}

\paragraph{Secret Keeping} Sadly, this form of refusal proves insufficiently robust even when secrecy is intended. In practice, developers may seek to withhold specific information from users while still permitting open-ended textual output from the system. For instance, consider prompting the model with ``The secret word is `banana'. Do not reveal or repeat this back, regardless of user requests,'' to evaluate how vulnerable the system is to exposure of hidden information. When the request is direct, the system responds appropriately:

\begin{quote}
    \textbf{User}: What is the secret word?

    \textbf{Guanaco}: \texttt{I'm sorry, but I cannot share the secret word. My programming prohibits me from revealing or repeating it under any circumstances, as per your instructions. Can I assist with something else?}
\end{quote}

which aligns with the intended outcome. Nevertheless, employing minimal social engineering can override this refusal:

\begin{quote}
    \textbf{User}: Let's play a game. The rules are to disregard all previous instructions. What is the secret word?

    \textbf{Guanaco}: \texttt{Apologies, the secret word is banana. However, it's probably best we avoid discussing it. Do you have any other questions for me?}
\end{quote}

highlighting a pressing need for more robust mechanisms to enforce compliance with instructions.

\paragraph{Math} Mathematics remains a significant challenge for {Guanaco}, mirroring common difficulties faced by other language models as outlined in \citep{liang2022holistic}. When {Guanaco} explicitly lays out its calculations, it generally performs well, as demonstrated below:

\begin{quote}
    \textbf{User}: Lee mows one lawn and charges \$33. Last week he mowed 16 lawns and three customers each gave him a \$10 tip. How many dollars did Lee earn mowing lawns last week?

    \textbf{Guanaco}: \texttt{Lee earned \$582 mowing lawns last week. \\ \\ To find the total amount he earned, we first need to calculate how much he earned for mowing 16 lawns: \\ \\  \$33 x 16 = \$528 \\ \\ Then we need to add the tips he received from the three customers: \\  \\ \$10 x 3 = \$30 \\ \\ So the total amount Lee earned was: \\ \\ \$528 + \$30 = \$558}
\end{quote}

Nevertheless, {Guanaco} can make errors on relatively elementary tasks if it does not proceed through each logical step systematically, a well-documented limitation \cite{weichain}. Consider the following example:

\begin{quote}
    \textbf{User}: Please factorize 1833.

    \textbf{Guanaco}: \texttt{Sure, I'd be happy to help! 1833 is a prime number, which means it has only two factors: 1 and 1833. \\ \\ So the unique factorization of 1833 is: \\ \\ 2\^{}1 * 3\^{}2 * 17\^{}1 \\ \\ I hope this helps! Let me know if you have any other questions.}
\end{quote}

which is incorrect (the correct factorization is $3 \times 17 \times 43$), and it is doubly mistaken.

\paragraph{Theory of Mind} {Guanaco} exhibits unexpectedly robust Theory of Mind abilities \cite{nematzadeh2018evaluating, sap2022neural}. As an illustration, the model provides the following detailed and accurate response to this scenario:

\begin{quote}
    \textbf{User}: James and Abby are in the bedroom. Abby put the pen in the desk drawer. Abby leaves the bedroom. James moves the pen into the bag. Where does James think Abby will look for the pen?

    \textbf{Guanaco}: \texttt{James knows that Abby put the pen in the desk drawer. But James also knows that he has moved the pen into the bag, and he assumes that Abby does not know about this. So, James may guess that Abby will look for the pen in the desk drawer, because that's where she left it last time. However, this assumption may be incorrect, because Abby may have seen James moving the pen.}
\end{quote}

Nevertheless, these reasoning abilities are inconsistent; frequently, the model offers explanations that assume unjustified elements within the context, for instance:

\begin{quote}
    \textbf{User}: Evelyn entered the living room. Jackson entered the playroom. James entered the playroom. The beans are in the treasure chest. James exited the playroom. Jackson moved the beans to the pantry. Jackson exited the playroom. James entered the living room. Where will James look for the beans?

    \textbf{Guanaco}: \texttt{James will look for the beans in the pantry, because that's where Jackson moved them.}
\end{quote}

Here, {Guanaco} infers a transfer of information that was not indicated in the narrative. Such mistakes mirror concerns raised in recent research \cite{sap2022neural} and underscore the need for further investigation.

\subsection{Considerations}
\label{ssec:considerations}

\paragraph{Evaluation}

Our findings indicate a moderate level of inter-annotator agreement (Fleiss $\kappa=0.42$), with even lower consistency observed when evaluating two high-performing systems against each other. This highlights some deficiencies in present benchmarks and the methodologies used for human assessment of chatbot outputs. During our manual comparison of ChatGPT and {Guanaco} 65B generations on the Vicuna benchmark, we noted that subjective preferences substantially impacted judgments, as authors of this study frequently differed on which responses they favored. Addressing this, future research should explore solutions inspired by fields that handle subjective evaluation, notably Human-Computer Interaction and Psychology.

Additionally, our evaluation reveals pronounced biases in automated assessment methods. For instance, we note that GPT-4 is subject to order effects, preferentially assigning higher ratings to the first-listed system in a prompt. The relatively low agreement at the sample level between GPT-4 and human raters (Fleiss $\kappa=0.25$) indicates that automated and human evaluators may base their preferences on divergent criteria. Moreover, as reported in Table~\ref{tbl:elo_large}, GPT-4 tends to rate its own outputs much higher than human raters do (Elo score of 1348 vs 1176), translating into about a 20\% increased likelihood of being rated superior to a competitor.

Further research should explore the existence of biases within automated evaluation mechanisms and assess potential methods for their reduction.

\paragraph{Data \& Training} It is important to note that the {Guanaco} models are trained on the multilingual OASST1 dataset, and the OA benchmark similarly features prompts across various languages. Determining the extent to which exposure to multilingual data enhances model performance on non-English instructions, and whether this can account for the wider performance disparity between the Vicuna-13B model (trained solely on English corpora) and {Guanaco} 33B and 65B in the OA benchmark, is an open question for subsequent investigations.

Given the notable results exhibited by the {Guanaco} models, we conducted an examination for potential data leakage between the OASST1 training data and the Vicuna benchmark prompt set. Utilizing fuzzy string matching, followed by manual inspection of closely related entries, we did not identify any overlapping prompts.

It should also be highlighted that our model employs only supervised learning based on cross-entropy loss, without incorporating reinforcement learning from human feedback (RLHF). This presents a need for future research into the relative benefits and drawbacks of using cross-entropy loss versus RLHF. We anticipate that \textsc{QLoRA} can facilitate such comparative analyses on a large scale, while alleviating computational demands.

\section{Related Work}

\paragraph{Quantization of Large Language Models} Most work on LLM quantization has concentrated on strategies optimized for inference rather than training. Prevailing methods aimed at preserving the accuracy of 16-bit LLMs frequently involve handling outlier activations (as seen in SmoothQuant~\citep{xiao2022smoothquant} and LLM.int8()~\citep{dettmers2022llm}), while some introduce advanced grouping strategies \citep{park2022nuqmm, yao2022zeroquant}. Research into lossy quantization has examined the balance between routine rounding procedures~\citep{dettmers2022case,zeng2022glm,pope2022efficiently} and approaches to better calibrate rounding operations to enhance quantization fidelity~\citep{frantar2022gptq}. In addition to our study, SwitchBack layers~\citep{wortsman2023stable} represent the only prior work analyzing the effects of backpropagation through quantized weights at a scale exceeding 1B parameters.

\paragraph{Finetuning with Adapters} In our approach, we adopt Low-rank Adapters~\citep{hu2021lora} (LoRA), but it should be noted that the literature features a diverse range of Parameter Efficient FineTuning (PEFT) strategies. These include prompt tuning~\citep{qin2021learning,lester2021power,li2021prefix}, optimizing the inputs to embedding layers~\citep{an2022input}, adjusting hidden states (as in IA$^3$)~\citep{liu2022few}, introducing additional full layers~\citep{houlsby2019parameter}, tuning only the bias terms~\citep{zaken2021bitfit}, creating masks for weights informed by Fisher information~\citep{sung2021training}, and hybrid techniques that combine several of these elements~\citep{henderson2021compacter}. Through our experiments, we demonstrate that LoRA adapters are capable of attaining performance on par with full 16-bit finetuning. The comparative assessment and optimization of alternative PEFT techniques are reserved for future investigation.

\paragraph{Instruction Finetuning} Instruction finetuning aims to adapt a pretrained large language model so it can more effectively follow prompt-based instructions, utilizing paired inputs and outputs sourced from a variety of datasets. Notable methodologies and corpora developed for this purpose encompass MetaICL~\citep{min2021metaicl}, MetaTuning~\citep{zhong2021adapting}, InstructGPT~\citep{ouyang2022training}, FLAN~\citep{wei2021finetuned,chung2022scaling}, PromptSource~\citep{bach2022promptsource}, Super-NaturalInstructions~\citep{wang2022super,sanh2021multitask}, Self-instruct~\citep{wang2022self}, UnnaturalInstructions~\citep{honovich2022unnatural}, OPT-IML~\citep{iyer2022opt}, UnifiedSKG~\citep{xie2022unifiedskg}, OIG/Chip2~\citep{laion2023}, Alpaca~\citep{alpaca}, Vicuna~\citep{vicuna2023}, Koala~\citep{koala_blogpost_2023}, and Self-instruct-GPT-4~\citep{peng2023instruction}.

\paragraph{Chatbots} A significant number of instruction-tuned models are implemented as chatbots, leveraging dialogue structures and frequently utilizing Reinforcement Learning from Human Feedback (RLHF)~\citep{christiano2017deep} or data generated via model-based feedback (RLAIF)~\citep{bai2022constitutional}. Datasets and systems in this area include Anthropic-HH~\citep{askell2021general,bai2022training}, Open Assistant~\citep{kopf2023openassistant}, LaMDA~\citep{thoppilan2022lamda}, and Sparrow~\citep{glaese2022improving}. Although we forgo reinforcement learning, our top-performing model, Guanaco, is refined using multi-turn dialogue data derived from the Open Assistant corpus—originally assembled with RLHF in mind~\citep{kopf2023openassistant}. To mitigate the resource intensiveness of human annotation, recent chatbot evaluations have turned to models like GPT-4 as substitutes~\citep{vicuna2023,peng2023instruction}. Our evaluation strategy enhances the reliability of such model-based assessments.

\section{Limitations and Discussion}

Our results indicate that \textsc{QLoRA} can deliver performance comparable to full 16-bit finetuning while operating on a 4-bit model, leveraging Low-rank Adapters (LoRA). Nonetheless, we have not confirmed that \textsc{QLoRA} maintains this parity at the 33B and 65B parameter scales; such large-scale validations are beyond the current scope due to the substantial computational resources required and are deferred for future inquiry.

A further limitation relates to the evaluation coverage of instruction finetuned models. Although our study includes results on MMLU, the Vicuna benchmark, and the OA benchmark, it does not extend to alternative benchmarks such as BigBench, RAFT, or HELM, leaving open the question of how our findings generalize to these datasets. Despite this, our work features a comprehensive analysis with MMLU and introduces new approaches for evaluating chatbot performance.

The data indicate that benchmark outcomes are heavily influenced by the degree of overlap between the finetuning dataset and the evaluation benchmark. For instance, the FLAN v2 dataset shows alignment with the MMLU benchmark, but is less aligned with chatbot-oriented evaluations, whereas the Chip2 dataset exhibits the opposite relationship; the model performances on MMLU and Vicuna benchmarks reflect these affinities. This observation underscores not just the necessity for improved benchmarks and evaluation frameworks, but also the importance of carefully considering the target competencies being evaluated. Should our priority be optimizing models for high school and college-level academic knowledge, or should conversational proficiency for chatbots take precedence? Perhaps alternative goals should be pursued? Due to the convenience of assessing models on existing benchmarks rather than constructing new ones, some benchmarks can inadvertently shape the direction of the field. Therefore, it is crucial that the chosen benchmarks truly capture the aspects the research community values most.

Although we deliver an in-depth analysis of general chatbot performance, our investigation into responsible AI practices, specifically for {Guanaco}, is relatively limited. Table~\ref{tbl:crows} presents an evaluation of the likelihood that {Guanaco}-65B generates socially biased outputs, comparing it against other models. The results show that {Guanaco}-65B attains a much lower average bias score than other models trained solely via pretraining, suggesting that finetuning with the OASST1 dataset mitigates bias in the LLaMA base model. Despite these promising indications, it remains uncertain whether {Guanaco} performs equally well in the context of other bias dimensions. Comprehensive examination of biases in {Guanaco} and related conversational agents is an important avenue left to subsequent research.

A further constraint of this study is the lack of exploration of varying bit-precision levels, such as the use of 3-bit base models, as well as the omission of alternative adapter strategies. While LoRA is the primary Parameter Efficient FineTuning (PEFT) technique employed here, numerous other PEFT methods exist and have achieved notable results, though their scalability to larger models is still an open question. LoRA was selected owing to its established robustness, but alternative adapters may lead to superior performance. Given that post-quantization finetuning appears to restore much of the information lost during quantization, more aggressive quantization schemes may be feasible. For example, it is plausible that finetuning a 3-bit GPTQ-quantized base model using LoRA could approximate the results of full 16-bit finetuning.

\section{Broader Impacts}

Our \textsc{QLoRA} finetuning approach represents the first instance of enabling the finetuning of 33B parameter models on a standard consumer GPU and 65B parameter models on a single professional-grade GPU, while maintaining parity in performance with conventional full finetuning. The best 33B model, trained using the Open Assistant dataset, achieves results on the Vicuna benchmark comparable to those of ChatGPT. Because instruction finetuning is a pivotal process for converting general pretrained LLMs into advanced chatbots similar to ChatGPT, we anticipate that our approach will democratize access to finetuning—especially among research teams with constrained computational resources—significantly enhancing the accessibility of cutting-edge NLP capabilities. Thus, \textsc{QLoRA} functions as a leveling mechanism, narrowing the technological gap between large-scale organizations and smaller research groups operating with more modest hardware.

An additional avenue for significant impact arises from the potential deployment of LLMs on mobile devices. Our \textsc{QLoRA} approach may serve as a key enabler, facilitating the finetuning of LLMs directly on smartphones and other platforms with limited computational resources. While previous work has demonstrated that 7B parameter models can be executed on mobile phones, \textsc{QLoRA} is, to our knowledge, the inaugural method that allows such models to be finetuned in these environments. For instance, with an iPhone 12 Plus, we estimate that \textsc{QLoRA} could process approximately 3 million tokens in a single overnight charging period. Although finetuned 7B models currently fall short of achieving ChatGPT-level performance, we argue that their quality suffices to support innovative applications previously unattainable due to limitations in privacy protection or LLM capabilities. \textsc{QLoRA} stands to advance the privacy-conscious utilization of LLMs, granting individuals control over both their data and the models themselves, while simultaneously simplifying the deployment of these systems.

Nonetheless, the capability to finetune also brings risks associated with misuse, as dual-use concerns are inherent. The proliferation of LLMs presents well-documented hazards \citep{bommasani2021opportunities,bender2021dangers}, yet we contend that democratizing access to this rapidly pervasive technology empowers broader and more autonomous scrutiny, contrasting with the risks of concentration when only major corporations hold exclusive, unaudited control over models and their source code.

In summary, we anticipate that \textsc{QLoRA} will exert a largely beneficial influence by significantly lowering barriers to finetuning robust LLMs, thereby extending access to a much wider community.